\section{Descripción del dataset}
\subsection{Datos y particiones}
Se utilizan los TFRecords de la competencia \textit{Petals to the Metal -- Flower
Classification on TPU} \citep{kagglePetalsMetal}, con 104 clases y fotografías
RGB. El conjunto incluye categorías relacionadas con la clasificación fina de
flores, estudiada también en Oxford Flowers \citep{nilsback2008}. Se mantienen
las particiones originales de Kaggle (Tabla~\ref{tab:particiones}).
\begin{table}[htbp]
\centering\small
\caption{Particiones originales utilizadas.}\label{tab:particiones}
\begin{tabular}{lrr}
\toprule Partición & Imágenes & Etiquetas locales \\
\midrule Train & 12,753 & Sí \\ Validación & 3,712 & Sí \\
Test & 7,382 & No \\ Total & 23,847 & \\
\bottomrule
\end{tabular}
\end{table}
Cada registro etiquetado contiene una imagen JPEG y una clase entera de 0 a 103;
test proporciona identificadores para generar predicciones. Se leen los archivos
de origen de $224\times224$, pero la entrada efectiva de todos los modelos
analizados es $112\times112\times3$. No se crea una partición nueva ni se usan
imágenes de validación para actualizar pesos.

\subsection{Desbalance y métricas}
El conteo histórico de train varía entre 18 y 782 imágenes por clase; su promedio
es 122.6. La exactitud global puede favorecer clases numerosas. Por ello el
criterio principal es F1 macro, que pondera por igual las 104 categorías.
Para una clase $c$, con verdaderos positivos $TP_c$, falsos positivos $FP_c$
y falsos negativos $FN_c$:
\begin{equation}
P_c=\frac{TP_c}{TP_c+FP_c},\qquad R_c=\frac{TP_c}{TP_c+FN_c},
\end{equation}
\begin{equation}
F1_c=\frac{2P_cR_c}{P_c+R_c},\qquad F1_{\rm macro}=\frac{1}{104}\sum_{c=0}^{103}F1_c.
\end{equation}
Los denominadores nulos se resuelven asignando cero a la métrica. El soporte es
el número de imágenes reales de esa clase en validación. Se informa también
exactitud y entropía cruzada: esta última penaliza errores de alta confianza
que un cambio de clase predicha no refleja por sí solo.
